% ==============================================================
\section{Related Work}\label{sec:related}
% ==============================================================

\subsection{Frequent Pattern Mining}

A comprehensive survey of frequent pattern mining is provided in~\cite{fournier2017survey}.
Apriori~\cite{agrawal1994fast} established the anti-monotone candidate generation and pruning framework,
while FP-Growth~\cite{han2000mining} and ECLAT~\cite{zaki2000scalable}
achieved speedups via FP-tree and vertical representations, respectively.
These methods target global frequency and do not model temporal variation.

\subsection{Temporal and Periodic Pattern Mining}

Segura-Delgado et al.~\cite{segura2020temporal} provide a comprehensive survey of temporal ARM.
PFPM~\cite{tanbeer2009discovering} and PPFPM~\cite{kiran2017discovering}
are periodicity mining methods targeting patterns that recur at fixed intervals,
while LPFIM~\cite{mahanta2005finding} detects frequent itemsets within time intervals
using gap-based thresholds.
LPPM~\cite{fournier2021mining} integrates locality and periodicity criteria.
These methods detect \emph{when and at what cycle} a pattern is frequent,
but do not model the \emph{cause} of variation.

\subsection{Apriori-window: Dense Interval Mining}

Apriori-window~\cite{dsaa2025} differs in design philosophy from the above methods.
It has no gap parameter or periodicity assumption, instead scanning a fixed-width
sliding window $W$ over a transaction sequence and directly evaluating
the local support $s_P(t)$ (Definition~\ref{def:support_series}) at each position $t$.
The output is the \emph{dense interval} $\mathcal{I}_P$ for each pattern $P$
(Definition~\ref{def:dense_interval})---the maximal contiguous range of window positions
where support continuously meets threshold $\theta$---representing
``during which period the pattern densely occurred.''
Candidate pattern space is efficiently pruned using Apriori's anti-monotonicity.
This work adopts Apriori-window as Step~0, using its output
$(\mathcal{F}, \{\mathcal{I}_P\})$ as input to the event attribution pipeline.

\subsection{Pattern Change Monitoring and Diagnostics}

Agrawal et al.~\cite{agrawal1995active} proposed trend detection via shape queries on rule statistics,
Chakrabarti et al.~\cite{chakrabarti1998mining} proposed a surprise metric based on description length,
and Liu et al.~\cite{liu2001fundamental} proposed extraction of ``fundamental changes'' between two time points.
These methods detect and characterize changes but lack a framework for attribution to external events.

\subsection{Change Point Detection}

Change point detection methods~\cite{aminikhanghahi2017survey}
represented by the CUSUM algorithm~\cite{page1954continuous}
identify structural changes in individual time series,
but are not equipped for the multiple testing burden of thousands of
pattern support series simultaneously or for linking changes to events.

\subsection{Event Attribution, Intervention Analysis, and Causal Inference}

Interrupted time series regression (ITS)~\cite{bernal2017interrupted},
the event study method~\cite{mackinlay1997event},
CausalImpact~\cite{brodersen2015causalimpact},
and the Wilcoxon rank-sum test~\cite{mann1947test}
test the effect of a known intervention on an outcome.
However, these methods typically target a single or small number of outcome variables
and cannot be directly applied to settings where thousands of pattern support series
are simultaneously tested against multiple events.

\subsection{Statistical Pattern Mining and Multiple Testing Correction}

Webb~\cite{webb2007discovering,webb2011filtered} addressed the problem of redundant significant patterns,
introducing self-sufficiency and productivity criteria to improve the relevance of discovered patterns.
Terada et al.~\cite{terada2013statistical} minimized the Bonferroni factor via LAMP,
and Llinares-L\'{o}pez et al.~\cite{llinares2015fast} achieved fast significant pattern mining
combining permutation testing with Tarone's testability criterion.
These methods test frequency differences between two groups, whereas this work tests
the association between \emph{temporal variation} of support and \emph{external events}.
Our pipeline adopts circular-shift permutation testing~\cite{davison1997bootstrap,yuan2024rigorous}
based on the surrogate data method~\cite{theiler1992testing}
and global BH correction~\cite{benjamini1995controlling,benjamini2001control}
to achieve FDR-controlled detection with higher power than Bonferroni correction~\cite{dunn1961multiple}.

\subsection{Summary of Differences}

Table~\ref{tab:related_comparison} summarizes whether each method family
can address the problem setting of this work.
To the best of our knowledge, this work is the first integrated framework that simultaneously attributes
many pattern time series to multiple events, controls multiple testing,
and eliminates secondary patterns.

\begin{table}[t]
\caption{Functional Comparison with Related Methods}
\label{tab:related_comparison}
\centering
\footnotesize
\begin{tabular}{lccccc}
\toprule
Method & \rotatebox{60}{TS Change} & \rotatebox{60}{Event Attr.} & \rotatebox{60}{Many Patterns} & \rotatebox{60}{Multi-test} & \rotatebox{60}{Dedup} \\
\midrule
Temporal FPM    & -- & -- & \checkmark & -- & -- \\
Apriori-window  & \checkmark & -- & \checkmark & -- & -- \\
Change Point    & \checkmark & -- & -- & -- & -- \\
ITS/Intervention & \checkmark & \checkmark & -- & -- & -- \\
Event Study     & \checkmark & \checkmark & -- & -- & -- \\
CausalImpact    & \checkmark & \checkmark & -- & -- & -- \\
Wilcoxon        & -- & \checkmark & -- & -- & -- \\
Statistical FPM & -- & -- & \checkmark & \checkmark & -- \\
\textbf{Proposed} & \checkmark & \checkmark & \checkmark & \checkmark & \checkmark \\
\bottomrule
\end{tabular}
\end{table}
